% Source: https://artofproblemsolving.com/community/c3h1270319_racing_and_barn_help
% Stable identity: aops-topic-1270319-q02
Put the corner where the rope is tied at the origin, with the barn occupying the rectangle $0\le x\le 50$, $0\le y\le 30$.

\begin{asy}
usepackage("amsmath");
size(6cm);
defaultpen(fontsize(10pt));
real r=60;
fill((0,0)--arc((0,0),r,90,360)--cycle,palegreen);
fill((50,0)--(60,0)--arc((50,0),10,0,90)--cycle,lightblue);
fill((0,30)--(0,60)--arc((0,30),30,90,0)--cycle,pink);
filldraw((0,0)--(50,0)--(50,30)--(0,30)--cycle,lightgray,black);
draw(arc((0,0),r,90,360)); draw(arc((50,0),10,0,90)); draw(arc((0,30),30,0,90));
dot((0,0)); label("rope 60",(-30,-4),S);
draw((0,0)--(-60*0.7071,-60*0.7071),dashed);
label("barn",(25,15)); label("50",(25,0),S); label("30",(0,15),W);
label("$r=10$",(60,8),E); label("$r=30$",(14,44));
label("\mbox{}",(min(currentpicture,user=true)+max(currentpicture,user=true))/2,black); // keeps the text after the figure black when LaTeX typesets the labels inline
\end{asy}

\begin{itemize}
\item Without bending around the barn the bull reaches three quarters of a circle of radius $60$: area $\dfrac{3}{4}\pi\cdot 60^2=2700\pi$.
\item Going along the $50$-foot wall and around its far corner, $10$ feet of rope remain, giving a quarter circle of radius $10$: area $25\pi$.
\item Going along the $30$-foot wall and around its far corner, $30$ feet of rope remain, giving a quarter circle of radius $30$: area $225\pi$.
\end{itemize}
The two small quarter circles lie behind different walls and do not overlap (one lies in $50\le x\le 60$, $0\le y\le10$, the other in $0\le x\le 30$, $30\le y\le 60$). The total area is
\[2700\pi+25\pi+225\pi=\boxed{2950\pi}\approx 9268\ \text{square feet}.\]
